\section{Data Collection}
\label{sec:data-collection}

\subsection{ASL Baseline Data}

The baseline dataset is the ``ASL Alphabet'' dataset on Kaggle~\cite{asl_alphabet_kaggle}, published under the GPL~2 license. It is downloaded with \texttt{data\_scripts/download\_data.py}, which reads the dataset identifier and the destination folder from \texttt{configs/data.yaml}. The images are never committed to the repository.

After download, the training folder contains 29 classes with 3{,}000 JPEG images each, 87{,}000 images in total: the letters A to Z and three extra classes stored as \texttt{del}, \texttt{nothing} and \texttt{space}. The separate test folder contains 28 images. These counts come from listing the downloaded files.

The dataset is American Sign Language. It is used only to build and test the data pipeline and as a comparison point; it is not PSL data.

\subsection{PSL Data Collected by the Team}

\todobox{\textbf{Owner: Aliza.} Describe the PSL collection and provide evidence:
\begin{itemize}[leftmargin=1.5em, topsep=2pt]
  \item The collection protocol: which letters (and which letter set, A--Z or Urdu letters), hand positions, lighting, backgrounds, camera, and the minimum number of images per letter per person.
  \item The consent process and how privacy is protected (anonymous IDs such as \texttt{p01}, no names stored, images kept out of Git).
  \item How the sign list was checked against a PSL source or a PSL user, with a citation.
  \item The number of volunteers and images actually collected (real counts only).
  \item Where the dataset is hosted and how to request access.
  \item Collection evidence figure: \texttt{reports/figures/psl\_sample\_grid.png} (a grid of example images, one per letter, from volunteers who agreed to have them shown).
\end{itemize}}

\safefigure{psl_sample_grid.png}{Example PSL images, one per letter.}{fig:psl-sample-grid}
